\documentclass[12pt,a4paper]{article}

% Encoding and fonts
\usepackage[utf8]{inputenc}
\usepackage[T1]{fontenc}
\usepackage{lmodern}

% Page layout
\usepackage[top=1in, bottom=1in, left=1in, right=1in]{geometry}

% Typography
\usepackage{microtype}
\usepackage{parskip}

% Language
\usepackage[english]{babel}

% Headers and footers
\usepackage{fancyhdr}
\pagestyle{fancy}
\fancyhf{}
\fancyhead[L]{\leftmark}
\fancyhead[R]{\thepage}
\fancyfoot[C]{\thepage}
\renewcommand{\headrulewidth}{0.4pt}
\renewcommand{\footrulewidth}{0pt}

% Section styling
\usepackage{titlesec}
\titleformat{\section}{\Large\bfseries}{\thesection}{1em}{}
\titleformat{\subsection}{\large\bfseries}{\thesubsection}{1em}{}
\titleformat{\subsubsection}{\normalsize\bfseries}{\thesubsubsection}{1em}{}
\titlespacing*{\section}{0pt}{2.5ex plus 1ex minus .2ex}{1.5ex plus .2ex}
\titlespacing*{\subsection}{0pt}{2ex plus 1ex minus .2ex}{1ex plus .2ex}
\titlespacing*{\subsubsection}{0pt}{1.5ex plus 1ex minus .2ex}{0.8ex plus .2ex}

% List formatting
\usepackage{enumitem}
\setlist[itemize]{topsep=0.5ex, itemsep=0.25ex, parsep=0pt}
\setlist[enumerate]{topsep=0.5ex, itemsep=0.25ex, parsep=0pt}

% Essential packages
\usepackage{graphicx}
\usepackage[table]{xcolor}
\usepackage{booktabs}
\usepackage{array}
\usepackage{tabularx}
\usepackage{longtable}
\usepackage{amsmath}
\usepackage{amssymb}
\usepackage[normalem]{ulem}
\rowcolors{2}{gray!10}{white}
\renewcommand{\arraystretch}{1.3}

% Cross-references
\usepackage{hyperref}
\usepackage{cleveref}

% Hyperref setup
\hypersetup{
    colorlinks=true,
    linkcolor=blue,
    filecolor=magenta,
    urlcolor=cyan,
}

% Code listings
\usepackage{listings}
\definecolor{codebg}{RGB}{248,248,248}
\definecolor{codeframe}{RGB}{200,200,200}
\definecolor{codekeyword}{RGB}{0,0,180}
\definecolor{codecomment}{RGB}{0,128,0}
\definecolor{codestring}{RGB}{163,21,21}
\lstset{
    basicstyle=\ttfamily\small,
    breaklines=true,
    breakatwhitespace=false,
    frame=single,
    framerule=0.5pt,
    rulecolor=\color{codeframe},
    backgroundcolor=\color{codebg},
    keepspaces=true,
    showstringspaces=false,
    tabsize=4,
    xleftmargin=1em,
    xrightmargin=1em,
    aboveskip=1em,
    belowskip=1em,
    keywordstyle=\color{codekeyword}\bfseries,
    commentstyle=\color{codecomment}\itshape,
    stringstyle=\color{codestring},
}

% YAML language definition for listings
\lstdefinelanguage{YAML}{
    keywords={true,false,null,y,n},
    keywordstyle=\color{blue}\bfseries,
    sensitive=false,
    comment=[l]{\#},
    morecomment=[s]{/*}{*/},
    commentstyle=\color{gray}\ttfamily,
    stringstyle=\color{red}\ttfamily,
    morestring=[b]',
    morestring=[b]",
}

% TOML language definition for listings
\lstdefinelanguage{TOML}{
    keywords={true,false},
    keywordstyle=\color{blue}\bfseries,
    sensitive=true,
    comment=[l]{\#},
    commentstyle=\color{gray}\ttfamily,
    stringstyle=\color{red}\ttfamily,
    morestring=[b]',
    morestring=[b]",
}

% Dockerfile language definition for listings
\lstdefinelanguage{Dockerfile}{
    keywords={FROM,RUN,CMD,LABEL,EXPOSE,ENV,ADD,COPY,ENTRYPOINT,VOLUME,USER,WORKDIR,ARG,ONBUILD,STOPSIGNAL,HEALTHCHECK,SHELL},
    keywordstyle=\color{blue}\bfseries,
    sensitive=false,
    comment=[l]{\#},
    commentstyle=\color{gray}\ttfamily,
    stringstyle=\color{red}\ttfamily,
    morestring=[b]',
    morestring=[b]",
}

% JSON language definition for listings
\lstdefinelanguage{JSON}{
    keywords={true,false,null},
    keywordstyle=\color{blue}\bfseries,
    sensitive=true,
    commentstyle=\color{gray}\ttfamily,
    stringstyle=\color{red}\ttfamily,
    morestring=[b]",
}

% Code listing float environment
\usepackage{float}
\newfloat{listing}{htbp}{lol}[section]
\floatname{listing}{Listing}

% Admonition boxes
\usepackage{tcolorbox}
\tcbuselibrary{skins,breakable}

% Admonition style definitions
\newtcolorbox{admonitionbox}[2][]{
    enhanced,
    breakable,
    colback=#2!5,
    colframe=#2!80!black,
    fonttitle=\bfseries,
    title=#1,
    left=2mm,
    right=2mm,
}

% Synthon tuple boxes
\newtcolorbox{synthonbox}{
    enhanced,
    colback=blue!5,
    colframe=blue!75!black,
    fontupper=\large,
    center upper,
    boxrule=1pt,
    arc=4pt,
}

\begin{document}

╔══════════════════════════════════════════════════════════════════════════════╗
║          SYNTHONICON ALGEBRAIC OPERATIONS --- TENSOR-MATH EDITION             ║
║  meet(⊓) $\cdot$ join(⊔) $\cdot$ tensor(\otimes) $\cdot$ lift $\cdot$ path $\cdot$ monad $\cdot$ decompositions     ║
╚══════════════════════════════════════════════════════════════════════════════╝

━━━━━━━━━━━━━━━━━━━━━━━━━━━━━━━━━━━━━━━━━━━━━━━━━━━━━━━━━━━━━━━━━━━━━━━━
§1  MEET  (⊓)  ---  Greatest Lower Bound
━━━━━━━━━━━━━━━━━━━━━━━━━━━━━━━━━━━━━━━━━━━━━━━━━━━━━━━━━━━━━━━━━━━━━━━━

────────────────────────────────────────────────────────────────────────
1.1  meet(Hv1\_human\_open, AtHv1\_primed)
────────────────────────────────────────────────────────────────────────

SETUP:
Hv1\_human\_open  : \langleD\_$\land$; T\_⋈; $R_{\supseteq}$; P\_+-; $F_{\hbar}$; K\_mod; G\_ב; $\Gamma$\_$\land$(SEL); $\Phi$\_c; $\Omega$\_0\rangle
AtHv1\_primed    : \langleD\_$\land$; T\_⋈; $R_{\supseteq}$; P\_+-; $F_{\hbar}$; K\_mod; G\_ב; $\Gamma$\_$\land$(SEL); $\Phi$\_sub; $\Omega$\_0\rangle

TENSOR-MATH ANALOGY:
In a product of bounded lattices L₁ $\times$ L₂ $\times$ ... $\times$ L\_n,
the meet is computed component-wise. Here $\Phi$-space is a 3-element
lattice  $\Phi$\_sub \textless{} ? \textless{} $\Phi$\_c  where $\Phi$\_c is the TOP element (absorbing).
So ⊓ on $\Phi$ returns $\Phi$\_c regardless of which operand holds it.
This is NOT the usual infimum --- it is an absorbing element,
analogous to the ''\top element in a co-Heyting algebra.''

{[}RESULT{]}
STATUS  : PASS
⟹  $\Phi$: Phi\_c ⊓ Phi\_sub $\rightarrow$ Phi\_c ($\Phi$\_c dominates in meet)

INSIGHT:
d(AtHv1\_primed, PsHv1\_constitutive) = 0.000 --- the gymnosperm
channel is algebraically identical to mechanically primed Arabidopsis.
The meet is the algebraic proof of cross-species functional conservation.

────────────────────────────────────────────────────────────────────────
1.2  meet(Hv1\_human\_open, 2GBI\_inhibitor) --- correct conflict
────────────────────────────────────────────────────────────────────────

SETUP:
Hv1\_human\_open  : T\_⋈  (H-bond network topology, cyclic)
2GBI\_inhibitor  : T\_$\in$  (condensed bicyclic, rigid ring network)
+ P\_$\pm$\^{}\psi (pseudosymmetric) vs P\_+- (directional)

TENSOR-MATH ANALOGY:
In module theory: two modules M, N have a non-trivial meet (intersection)
only when they are sub-objects of a common ambient module.
T\_⋈ and T\_$\in$ are NOT in the same T-equivalence class --- there is no
common parent topology. The meet is \perp: undefined.

{[}RESULT{]}
STATUS  : CONFLICT
  T
  P
⟹  $\Phi$: Phi\_c ⊓ Phi\_sub $\rightarrow$ Phi\_c ($\Phi$\_c dominates in meet)
⟹  F: F\_hbar ⊓ F\_eth $\rightarrow$ F\_eth (min)

INSIGHT:
The conflict is the answer. 2GBI occludes the channel --- it does NOT
adopt the channel's topology or merge with it. The drug sits in the
pore and blocks it physically. tensor() is the correct operation
for that question (see §3). meet() correctly returns \perp.

────────────────────────────────────────────────────────────────────────
1.3  Topological protection order: meet(Cooper pair, TI surface)
────────────────────────────────────────────────────────────────────────

SETUP:
cooper\_pair           : $\Omega$\_Z   (ℤ winding number; BCS s-wave)
topological\_insulator : $\Omega$\_Z₂  (ℤ₂ time-reversal; Bi₂Se₃ class)

$\Omega$ ORDINAL:  TRIVIAL(0) \textless{} Z₂(1) \textless{} Z(2) \textless{} CHERN(3) \textless{} NON\_ABELIAN(4)

TENSOR-MATH ANALOGY:
The $\Omega$ lattice mirrors the Altland-Zirnbauer (AZ) classification.
Under meet, $\Omega$ behaves like a meet-semilattice: the conservative
guarantee is the WEAKER protection class --- you can guarantee the
physics at the intersection only up to the least protected invariant.
This is the lattice-theoretic reason topological surface states
are fragile to symmetry-breaking perturbations that act on the
weaker Z₂ invariant.

NOTE: This is a topology-focused meet. D and T will conflict here
(MOLECULAR vs SUPRAMOLECULAR, BOWTIE vs BOWTIE --- actually BOWTIE matches).
The pedagogical point is the $\Omega$ ordering rule.

{[}RESULT{]}
STATUS  : CONFLICT
  D
  P
  $\Gamma$
⟹  $\Phi$: Phi\_c ⊓ Phi\_sub $\rightarrow$ Phi\_c ($\Phi$\_c dominates in meet)
⟹  $\Omega$: Omega\_Z ⊓ Omega\_Z2 $\rightarrow$ Omega\_Z2 (min protection)
⟹  K: K\_slow ⊓ K\_fast $\rightarrow$ K\_slow (min)

━━━━━━━━━━━━━━━━━━━━━━━━━━━━━━━━━━━━━━━━━━━━━━━━━━━━━━━━━━━━━━━━━━━━━━━━
§2  JOIN  (⊔)  ---  Least Upper Bound  /  Design Target
━━━━━━━━━━━━━━━━━━━━━━━━━━━━━━━━━━━━━━━━━━━━━━━━━━━━━━━━━━━━━━━━━━━━━━━━

────────────────────────────────────────────────────────────────────────
2.1  join(Hv1\_human\_open, PsHv1\_constitutive)
────────────────────────────────────────────────────────────────────────

SETUP:
Hv1\_human\_open      : $\Phi$\_c  (H-bond water chain)
PsHv1\_constitutive  : $\Phi$\_sub (constitutively primed; no $\Phi$\_c required)

TENSOR-MATH ANALOGY:
In a module category, join is the pushout: the minimal object
receiving morphisms from both M and N. Here the join forces $\Phi$\_c
into the design target --- the coproduct must accommodate the more
demanding constraint (criticality) from the human channel.

{[}RESULT{]}
STATUS  : PASS
⟹  $\Phi$: join propagates Phi\_c (criticality is join-dominant)

────────────────────────────────────────────────────────────────────────
2.2  join(2GBI\_inhibitor, HIF\_inhibitor) --- no common design target
────────────────────────────────────────────────────────────────────────

SETUP:
2GBI : T\_$\in$  (condensed fused bicyclic: charge-delocalised ring system)
HIF  : T\_\textbar{}  (flexible linear linker: two pharmacophores, conformational freedom)

TENSOR-MATH ANALOGY:
These are objects from DIFFERENT categories: T\_$\in$ is in the ''network''
homotopy class; T\_\textbar{} is in the ''linear'' class. The coproduct (join)
requires a common ambient object --- but no scaffold topology subsumes
both rigid bicyclic AND flexible linear in the same T class.
The join is undefined (\perp in the join-semilattice for T).

DESIGN CONSEQUENCE:
This algebraically forbids ''best of both worlds'' scaffold merging.
The correct design path is tensor() (§3.3): what do they predict
TOGETHER as a co-assembly, not as a single molecule?

{[}RESULT{]}
STATUS  : CONFLICT
  T

────────────────────────────────────────────────────────────────────────
2.3  join(imatinib, GNF2) --- combined Type II + allosteric target
────────────────────────────────────────────────────────────────────────

SETUP:
imatinib  : K\_slow, G\_local,  $\Gamma$\_$\land$(SPECIFIC), $\Phi$\_sub
GNF-2     : K\_mod,  G\_gimel,  $\Gamma$\_$\rightarrow$(SELECTIVE), $\Phi$\_c

TENSOR-MATH ANALOGY:
join on ordered dimensions takes supremum: max(K\_slow, K\_mod) = K\_slow
(the slowest kinetics must be satisfied by the design target).
G: max(G\_ב, G\_ג) = G\_ג (target must have mesoscale propagation).
$\Phi$\_c propagates. The join is the ''hardest'' requirement from either
source --- a design constraint satisfaction problem whose solution
is the join element.

Note: T conflict expected (T\_⋈ vs T\_branched). This is the real
bottleneck: combining the DFG-out loop topology with the myristoyl
pocket branched topology requires T-class redesign.

{[}RESULT{]}
STATUS  : CONFLICT
  T
  $\Gamma$
⟹  $\Phi$: join propagates Phi\_c (criticality is join-dominant)
⟹  F: F\_hbar ⊔ F\_eth $\rightarrow$ F\_hbar (max)
⟹  K: K\_slow ⊔ K\_mod $\rightarrow$ K\_mod (max)
⟹  G: G\_beth ⊔ G\_gimel $\rightarrow$ G\_gimel (max)

INSIGHT:
The T-conflict reveals the hard part of dual-mechanism ABL inhibitor
design: the DFG-out pocket (T\_⋈, cyclic coordination) and the
myristoyl pocket (T\_branched, pendant ligand) cannot be joined
without a new T topology. This correctly predicts why Type II +
allosteric dual-binders have been rare and difficult to design.

━━━━━━━━━━━━━━━━━━━━━━━━━━━━━━━━━━━━━━━━━━━━━━━━━━━━━━━━━━━━━━━━━━━━━━━━
§3  TENSOR  (\otimes)  ---  Bifunctor / Co-Assembly Prediction
━━━━━━━━━━━━━━━━━━━━━━━━━━━━━━━━━━━━━━━━━━━━━━━━━━━━━━━━━━━━━━━━━━━━━━━━

────────────────────────────────────────────────────────────────────────
3.1  tensor(electron, hole) --- exciton precursor; bound state gap
────────────────────────────────────────────────────────────────────────

SETUP:
electron : D\_mol, T\_\textbar{}, P\_donor,   F\_high, K\_fast, G\_local, $\Omega$\_0
hole     : D\_mol, T\_\textbar{}, P\_acceptor, F\_high, K\_fast, G\_local, $\Omega$\_0

KEY PREDICTION: T\_\textbar{} \otimes T\_\textbar{} $\rightarrow$ T\_\textbar{}  (no topology promotion for SAME input)

TENSOR-MATH ANALOGY:
In Hilbert space: ℋ\_e \otimes ℋ\_h gives the two-particle space.
The Coulomb interaction H\_Coulomb is an element of End(ℋ\_e \otimes ℋ\_h) ---
it is NOT included in the tensor product itself. The bound exciton
lives in a subspace selected by H\_Coulomb, which requires a separate
'meet with binding potential' operation in synthon algebra.

This example teaches the semantic boundary:
tensor()  = statistical co-occupancy (two-particle Hilbert space)
meet()    = intersection (bound-state subspace selection)
The exciton  = tensor(e, h) \textgreater{}\textgreater{} meet(coulomb\_binding\_potential\_synthon)

{[}RESULT{]}
STATUS  : PASS
⟹  P: P\_minus \otimes P\_plus $\rightarrow$ P\_pm\_pseudo (charge pairing)
⟹  \xi\_CP: 7.741 + 7.741 $-$ 0.5$\times$6.635 = 12.164 nats (6/7 primitives match $\rightarrow$ I $\approx$ 6.635 nat)
\xi\_CP    : 12.164 nats

INSIGHT (P$\rightarrow$DONOR\_ACCEPTOR):
P: DONOR \otimes ACCEPTOR $\rightarrow$ DONOR\_ACCEPTOR is the primordial signature of
charge-transfer: electron and hole have opposing charge asymmetry and
the tensor correctly predicts a directed dipole will form.
Frenkel exciton (T stays LINEAR) vs Wannier-Mott (T\_bowtie via binding)
distinguished by whether the binding step is allowed.

────────────────────────────────────────────────────────────────────────
3.2  tensor(phonon\_acoustic, magnon, $\lambda$=0.4) --- magnetoelastic polaron
────────────────────────────────────────────────────────────────────────

SETUP:
phonon\_acoustic : D\_supra, T\_\textbar{}, P\_sym,  F\_low, K\_fast, G\_global, $\Omega$\_0
magnon          : D\_supra, T\_\textbar{}, P\_sym,  F\_med, K\_mod,  G\_meso,   $\Omega$\_0

COUPLING: $\lambda$ = 0.4 (moderate spin-phonon coupling, e.g. YIG or MnF₂)

TENSOR-MATH ANALOGY:
The magnon-phonon Hamiltonian H\_mp = $\Sigma$ g\_kq (a\_k + a†\_k)(b\_q + b†\_q)
is a bilinear coupling in the product Fock space. In the synthon
algebra this is exactly tensor(): the co-assembly of two SUPRA
collective modes. The $\lambda$ parameter encodes g\_kq coupling strength ---
high $\lambda$ reduces \xi\_tensor via mutual information discount.

EXPECTED: G $\rightarrow$ G\_global (phonon bath extends over whole crystal);
F $\rightarrow$ F\_low (phonon decoherence is the bottleneck);
K $\rightarrow$ K\_fast (acoustic phonon velocity dominates).

{[}RESULT{]}
STATUS  : PASS
⟹  F: min(F\_ell, F\_eth) $\rightarrow$ F\_ell (bottleneck)
⟹  K: min(K\_fast, K\_mod) $\rightarrow$ K\_mod (trap risk propagates)
⟹  G: max(G\_aleph, G\_gimel) $\rightarrow$ G\_aleph (coarsest scale dominates)
⟹  $\Gamma$: Gamma\_or \otimes Gamma\_and $\rightarrow$ $\Gamma$\_$\land$(BROAD) (AND-composition)
⟹  \xi\_CP: 7.507 + 7.590 $-$ 0.4$\times$4.290 = 13.381 nats (4/7 primitives match $\rightarrow$ I $\approx$ 4.290 nat)
\xi\_CP    : 13.381 nats

────────────────────────────────────────────────────────────────────────
3.3  tensor(Majorana, Majorana, $\lambda$=0.3) --- topological qubit
────────────────────────────────────────────────────────────────────────

SETUP:
majorana\_fermion : D\_mol, T\_braid, P\_sym(self-conjugate), F\_high,
K\_trap, G\_local, $\Phi$\_c, $\Omega$\_NA

SPECIAL RULE: T\_braid \otimes T\_braid $\rightarrow$ T\_braid
(braided exchange statistics are preserved in co-assembly;
anyonic topology does NOT network-promote like spatial structures)

TENSOR-MATH ANALOGY:
The braid group B\_n has a natural tensor structure: B\_n \otimes B\_m \subseteq B\_\{n+m\}
as a sub-group. The braid topology of two Majorana modes tensors to
give a larger braid system, not a ''network'' of modes.
This is why T\_braid has its own special tensor rule: it does not
obey the standard topology promotion hierarchy.

$\Omega$ rule: $\Omega$\_NA \otimes $\Omega$\_NA $\rightarrow$ $\Omega$\_NA  (non-Abelian protection preserved)

PHYSICAL MEANING:
Two spatially separated Majorana modes in a Kitaev chain form a
topological qubit with $\Omega$\_NA protection. The tensor product predicts
this ensemble correctly: T\_braid preserved, $\Omega$\_NA preserved,
G: max(LOCAL, LOCAL) = LOCAL (both modes are localised).

{[}RESULT{]}
STATUS  : PASS
⟹  T: T\_braid \otimes T\_braid $\rightarrow$ T\_braid (anyonic statistics preserved under composition)
⟹  \xi\_CP: 8.314 + 8.314 $-$ 0.3$\times$8.314 = 14.134 nats (7/7 primitives match $\rightarrow$ I $\approx$ 8.314 nat)
\xi\_CP    : 14.134 nats

CONTRAST WITH EXAMPLE 3.1:
electron \otimes hole: T\_\textbar{} \otimes T\_\textbar{} $\rightarrow$ T\_\textbar{} (same topology, no promotion)
majorana \otimes majorana: T\_braid \otimes T\_braid $\rightarrow$ T\_braid (braid preserved)
magnon \otimes phonon: T\_\textbar{} \otimes T\_\textbar{} $\rightarrow$ T\_\textbar{} (same topology, no promotion)

\begin{lstlisting}
The ONLY cases where tensor changes T are cross-topology products:
T_| \otimes T_in -> T_in  (linear + network -> network)
T_in \otimes T_□ -> T_□  (network + cage -> cage)
This mirrors the intuition: co-assembly of a network and a linear
entity produces network-class organisation, not vice versa.
\end{lstlisting}

━━━━━━━━━━━━━━━━━━━━━━━━━━━━━━━━━━━━━━━━━━━━━━━━━━━━━━━━━━━━━━━━━━━━━━━━
§4  LIFT  ---  Natural Transformations Between Domain Categories
━━━━━━━━━━━━━━━━━━━━━━━━━━━━━━━━━━━━━━━━━━━━━━━━━━━━━━━━━━━━━━━━━━━━━━━━

────────────────────────────────────────────────────────────────────────
4.1  lift\_to\_temporal --- molecule enters catalytic cycle category
────────────────────────────────────────────────────────────────────────

RULES:
D\_$\land$ $\rightarrow$ D\_$\infty$          (point object $\rightarrow$ temporal / periodic)
T   $\rightarrow$ T\_⋈          (any $\rightarrow$ bowtie: substrate-in / product-out loop)
$R_{\supseteq}$ $\rightarrow$ R\_‡          (non-covalent $\rightarrow$ dynamic catalytic / turnover)
K\_fast $\rightarrow$ K\_mod     (fast-exchange adjusted for catalytic dwell time)
$\Gamma$   $\rightarrow$ $\Gamma$\_$\rightarrow$(SEL)     (sequential: binding-order enforced)

TENSOR-MATH ANALOGY:
In category theory: D\_$\infty$ objects are ''internal monoids'' --- objects
equipped with a self-map (the catalytic step). The temporal lift is
the functor that sends a morphism (binding event) to a monoid action
(the catalytic cycle). It is the ''loop-space'' functor $\Omega$ in topology:
$\Omega$(X) maps a pointed space to its loop space, endowing it with a
group structure. Here the ''group'' is the catalytic turnover.

COST: 0.0 nats (no thermodynamic cost to describe the catalytic frame)
GROUNDING REQUIRED: Axiom 6 --- must specify a reset event (product
release, cofactor regeneration). lift alone does not ground; grounding
is a separate obligation enforced by the Axiom validator.

Found temporal synthon in catalog: global\_supply\_chain
Notation: \langleD\_infinity; T\_network{[}1:\emph{{]}; R\_dagger; P\_minus; F\_eth; K\_mod; G\_aleph; Gamma\_and(SELECTIVE); Phi\_sub; 1:}\rangle

────────────────────────────────────────────────────────────────────────
4.2  lift\_to\_spatial --- molecule $\rightarrow$ crystal secondary building unit
────────────────────────────────────────────────────────────────────────

RULES:
D\_$\land$  $\rightarrow$ D\_△         (molecular $\rightarrow$ supramolecular crystal)
T    $\rightarrow$ T\_□          (any $\rightarrow$ hub-node: SBU topology)
G\_ב  $\rightarrow$ G\_ג (min)   (local $\rightarrow$ mesoscale; crystal requires mesoscale)
$\Gamma$    $\rightarrow$ $\Gamma$\_$\land$(SEL)    (coordinated multi-dentate recognition)

TENSOR-MATH ANALOGY:
The spatial lift is the ''classifying space'' functor B:
B(G) takes a group G (the molecule's local symmetry) and produces
a space whose loop space is G. In crystal engineering terms:
the SBU is the ''classifying object'' for the crystal packing symmetry.
The T\_□ (hub) topology encodes the branching valence of the SBU ---
how many struts radiate from the node.

A PRACTICAL CONSEQUENCE:
lift\_to\_spatial(2GBI\_inhibitor) would give T\_□ --- but 2GBI is
bicyclic (T\_$\in$). The T change ($\in$ $\rightarrow$ □) has a real cost in synthesis:
you must redesign the ring to have branching coordination points.
This is not a free transformation.

────────────────────────────────────────────────────────────────────────
4.3  criticality\_lift --- non-zero cost functor, fidelity gate
────────────────────────────────────────────────────────────────────────

GATE:  F $\geq$ $F_{\hbar}$ required to lift $\Phi$\_sub $\rightarrow$ $\Phi$\_c
COST:  +2.303 nats (one decade of probability: log\_e(10))

This is the primitive-space analog of the LANDAUER BOUND:
In information theory: erasing 1 bit costs k\_B$\cdot$T$\cdot$ln(2) = 0.693 nats.
Here: acquiring criticality (a binary phase) costs 2.303 nats.
The factor difference (2.303/0.693 $\approx$ 3.32) reflects the multi-dimensional
nature of a criticality transition vs a single-bit erasure.

TENSOR-MATH ANALOGY:
criticality\_lift is the functor from the ''generic'' category of
subcritical systems to the ''$\Phi$\_c-structured'' category where every
object has an associated RG fixed point. The cost 2.303 is the
''action'' of the functor --- the price of the structural promotion.

ASYMMETRY:
There is no ''criticality\_lower'' functor. Once $\Phi$\_c is in context.
the F-floor ratchet prevents downstream operations from reducing F.
This encodes the thermodynamic irreversibility of phase transitions:
you cannot un-boil an egg by running the monad backwards.

DEMO: lift on Hv1\_human\_closed (F\_high, $\Phi$\_sub) --- should PASS with cost.

LIFT BLOCKED: F gate not met for Hv1\_human\_closed
{[}BLOCKED{]}  lift(critical)  --- $\Phi$\_c lift not applicable: D\_$\infty$ or G $\geq$ G\_ג required for $\Phi$\_c eligibility

CONTRAST: What if F is LOW? (Try on a low-fidelity synthon)
conductingelectron (F\_high) $\rightarrow$ lift passes  (gate: F $\geq$ $F_{\hbar}$ met)
phonon\_acoustic   (F\_low)  $\rightarrow$ lift BLOCKED  (F\_low \textless{} $F_{\hbar}$ required)
This correctly encodes: phonons cannot be criticality-lifted; they are
not the order parameter. Only high-fidelity recognition modes can
undergo a symmetry-breaking transition.

━━━━━━━━━━━━━━━━━━━━━━━━━━━━━━━━━━━━━━━━━━━━━━━━━━━━━━━━━━━━━━━━━━━━━━━━
§5  PATH  ---  Geodesic in HotSwap Kleisli Category
━━━━━━━━━━━━━━━━━━━━━━━━━━━━━━━━━━━━━━━━━━━━━━━━━━━━━━━━━━━━━━━━━━━━━━━━

────────────────────────────────────────────────────────────────────────
5.1  path(AtHv1\_silent, AtHv1\_primed) --- discrete topology change
────────────────────────────────────────────────────────────────────────

SETUP:
AtHv1\_silent : T\_$\in$ (network; RSN locks topology)
AtHv1\_primed : T\_⋈ (bowtie; RSN removed, voltage-responsive loop)

HotSwap hard constraint: T must be identical along the path.
T\_$\in$ $\neq$ T\_⋈ $\rightarrow$ NO PATH in the HotSwap graph.

TENSOR-MATH ANALOGY:
In differential geometry: path-connectedness within a homotopy class.
T\_$\in$ and T\_⋈ are in DIFFERENT homotopy classes of the topology space.
There is no continuous deformation from T\_$\in$ to T\_⋈ --- the transition
requires a topological jump (equivalent to tearing the manifold).
This is a ''1st-order-like'' morphism: latent cost \textgreater{} 0 with no
intermediate states.

PHYSICAL MEANING:
Mechanical priming (membrane stretch) is not a smooth conformational
change. It releases the RSN kinetic trap all at once --- three primitives
(K, T, P) change simultaneously. d = 3.3 nats. No smooth path.

PATH BLOCKED: no path from AtHv1\_silent to AtHv1\_primed
Reason: T-class boundary (T\_network $\neq$ T\_bowtie)
This confirms: mechanical priming is a discrete 1st-order jump.

────────────────────────────────────────────────────────────────────────
5.2  path(2GBI\_inhibitor, HIF\_inhibitor) --- scaffold evolution barrier
────────────────────────────────────────────────────────────────────────

SETUP:
2GBI : T\_$\in$ (condensed bicyclic network; charge delocalised ring)
HIF  : T\_\textbar{} (flexible linear linker; two independent pharmacophores)

PATH STATUS: BLOCKED (T\_$\in$ $\neq$ T\_\textbar{})

DESIGN CONSEQUENCE:
Scaffold optimisation by chemical synthesis (adding substituents,
adjusting ring substitution) cannot bridge 2GBI $\rightarrow$ HIF.
The evolution from 2GBI-class to HIF-class is a DESIGN DISCONTINUITY:
it requires a new synthesis strategy, not incremental SAR.
This is why Papers 1$\rightarrow$2 (in the Webster/Tombola series) represent
a genuine paradigm shift in Hv1 inhibitor design, not iteration.

TENSOR-MATH ANALOGY:
In algebraic K-theory: the ''suspension'' isomorphism connects
K₀(T\_$\in$-class) to K₀(T\_\textbar{}-class) --- but only via an explicit
stabilisation construction (adding a trivial factor), which
corresponds to adding a flexible spacer and rebuilding the pharmacophore.
That IS the HIF design.

PATH BLOCKED: T\_$\in$ $\rightarrow$ T\_\textbar{} requires T-class redesign.
d(2GBI, HIF) = 1.500 nats
The distance encodes the design effort; the blocked path encodes irreversibility.

────────────────────────────────────────────────────────────────────────
5.3  path(topological\_insulator, conducting\_electron) --- topological gap
────────────────────────────────────────────────────────────────────────

SETUP:
topological\_insulator : D\_supra, T\_⋈, $\Omega$\_Z₂, $\Phi$\_sub
conducting\_electron   : D\_mol,   T\_\textbar{},  $\Omega$\_0,  $\Phi$\_sub

TWO BARRIERS:
1. D: SUPRAMOLECULAR $\neq$ MOLECULAR (bulk TI vs point particle)
2. T: T\_⋈ $\neq$ T\_\textbar{} (Dirac cone topology vs linear dispersion)

PATH STATUS: BLOCKED (two independent class barriers)

ASYMMETRY (SYNTHONICON\_LANG.md §3e):
path(TI $\rightarrow$ Fermi liquid) is blocked by D/T mismatch.
path(Fermi liquid $\rightarrow$ TI) is also blocked (reverse).
But the COST is asymmetric: TI $\rightarrow$ trivial metal costs $\Delta$\xi \textgreater{} 0
(destroying topological order releases entropy); trivial $\rightarrow$ TI costs
MORE (ordering requires investment). The F-floor ratchet encodes this:
once in the TI class (F\_high), the floor is raised and subsequent
operations cannot lower F.

PHYSICAL MEANING:
Topological protection IS the blocked path. A Majorana edge mode
cannot continuously deform to a trivial fermion without closing
the bulk gap. The HotSwap path-block is the algebraic encoding of
the bulk-boundary correspondence: the gapless surface is protected
precisely because there is no smooth path to the trivial phase.

PATH BLOCKED: topological gap is algebraically protected.
D-barrier: Dimensionality.SUPRAMOLECULAR $\neq$ Dimensionality.MOLECULAR
T-barrier: Topology.CYCLIC\_BOWTIE $\neq$ Topology.LINEAR
$\Omega$-gap:     TopoIndex.Z2\_CLASS $\rightarrow$ TopoIndex.TRIVIAL

━━━━━━━━━━━━━━━━━━━━━━━━━━━━━━━━━━━━━━━━━━━━━━━━━━━━━━━━━━━━━━━━━━━━━━━━
§6  PIPELINES  ---  Monadic Composition (SynthonM)
━━━━━━━━━━━━━━━━━━━━━━━━━━━━━━━━━━━━━━━━━━━━━━━━━━━━━━━━━━━━━━━━━━━━━━━━

────────────────────────────────────────────────────────────────────────
6.1  Hv1 cross-species conservation (from hv1\_paper\_reproduction.syn)
────────────────────────────────────────────────────────────────────────

PIPELINE:
start: Hv1\_human\_open
\textgreater{}\textgreater{}= meet(AtHv1\_primed)          \# both T\_⋈ --- no conflict; $\Phi$\_c dominates
\textgreater{}\textgreater{}= join(PsHv1\_constitutive)    \# d=0.000; near-trivial join
\textgreater{}\textgreater{}= assert($\Phi$ == $\Phi$\_c)            \# H-bond chain preserved
\textgreater{}\textgreater{}= assert(phi\_c\_score $\geq$ 0.35)  \# Varma weight for MOLECULAR/LOCAL

MONAD SEMANTICS:
Each \textgreater{}\textgreater{}= threads the current synthon through the next operation.
MaybeT: if any step returns None (conflict), the rest short-circuit.
WriterT: $\Delta$\xi accumulates --- here 0.0 at every step (all within class).
StateT: f\_floor set by join to F\_eth; criticality\_ok unchanged (no lift).

Result   : PASS
Output   : \langleD\_wedge; T\_bowtie; R\_superset; P\_directional; F\_hbar; K\_mod; G\_beth; Gamma\_and(SELECTIVE); Phi\_c\rangle
$\Delta$\xi total : 0.000 nats
F-floor  : F\_hbar

RESULT: join(meet(Hv1\_human\_open, AtHv1\_primed), PsHv1\_constitutive)
$\Delta$\xi\_CP = 0.000. Cross-species conservation demonstrated algebraically.

────────────────────────────────────────────────────────────────────────
6.2  mplus (\textless{}\textbar{}\textgreater{}) --- fallback design strategy
────────────────────────────────────────────────────────────────────────

PIPELINE LOGIC:
strategy\_A: start(Hv1\_open) \textgreater{}\textgreater{}= meet(2GBI)    \# will BLOCK (T-conflict)
strategy\_B: start(Hv1\_open) \textgreater{}\textgreater{}= meet(AtHv1\_primed)  \# will PASS

\begin{lstlisting}
result = strategy_A.mplus(strategy_B)
-> try A; on BLOCK, automatically fall back to B.
\end{lstlisting}

MONAD SEMANTICS (mplus = MonadPlus operation \textless{}\textbar{}\textgreater{}):
mplus is the ''choice'' combinator in the Maybe monad:
Nothing \textless{}\textbar{}\textgreater{} Just x = Just x
Here: BLOCKED \textless{}\textbar{}\textgreater{} PASS = PASS
Cost of the successful branch accumulates; failed branch cost is lost
(WriterT append-only; but MaybeT discards the failed branch output).

THIS IS ALGEBRAICALLY CORRECT:
mplus models branching design paths. In retrosynthesis: ''try this
route; if blocked by a protecting group conflict, try the alternative.''
The F-floor from the failed branch does NOT transfer (the failed
branch never raised the floor).

strategy\_A (meet 2GBI)    : BLOCKED
strategy\_B (meet AtHv1)   : PASS
mplus result              : PASS --- strategy\_B succeeded
Output : \langleD\_wedge; T\_bowtie; R\_superset; P\_directional; F\_hbar; K\_mod; G\_beth; Gamma\_and(SELECTIVE); Phi\_c\rangle

────────────────────────────────────────────────────────────────────────
6.3  Quantum co-assembly pipeline: Cooper pair \otimes TI surface
────────────────────────────────────────────────────────────────────────

PIPELINE:
start: cooper\_pair  ($\Omega$\_Z, $\Phi$\_c, T\_⋈)
\textgreater{}\textgreater{}= tensor(topological\_insulator, $\lambda$=0.3)
\textgreater{}\textgreater{}= assert(criticality\_phase == Phi\_c)

PHYSICAL MEANING:
Proximity effect: a Cooper pair condensate adjacent to a TI surface
can induce topological superconductivity. The tensor product predicts
the primitive structure of this co-assembly. The assert checks whether
$\Phi$\_c is preserved in the heterostructure.

$\Omega$ in tensor: max($\Omega$\_Z, $\Omega$\_Z₂) = $\Omega$\_Z (Z-class \textgreater{} Z₂-class in protection ordinal)
T: T\_⋈ \otimes T\_⋈ $\rightarrow$ T\_⋈ (bowtie preserved; Dirac + pairing loop)
$\Phi$: $\Phi$\_c \otimes $\Phi$\_sub $\rightarrow$ $\Phi$\_c (criticality propagates)

This tensor product is the algebraic encoding of the proximity effect
that generates topological superconductor phases (class D/DIII).

Result: PASS
Output : \langleD\_triangle; T\_bowtie; R\_superset; P\_pm\_pseudo; F\_hbar; K\_slow; G\_gimel; Gamma\_and(SELECTIVE); Phi\_c\rangle
$\Delta$\xi\_CP  : 13.854 nats
$\Omega$      : None  (higher $\Omega$ propagates in tensor)
$\Phi$      : CriticalityPhase.CRITICAL  ($\Phi$\_c propagates)

━━━━━━━━━━━━━━━━━━━━━━━━━━━━━━━━━━━━━━━━━━━━━━━━━━━━━━━━━━━━━━━━━━━━━━━━
§7  DECOMPOSITIONS  ---  Factor $\cdot$ Cofactor $\cdot$ Kernel $\cdot$ Principal $\cdot$ Project $\cdot$ Complement
━━━━━━━━━━━━━━━━━━━━━━━━━━━━━━━━━━━━━━━━━━━━━━━━━━━━━━━━━━━━━━━━━━━━━━━━

────────────────────────────────────────────────────────────────────────
7.1  cofactor(cooper\_pair, conducting\_electron) --- inverse tensor: what is the pairing partner?
────────────────────────────────────────────────────────────────────────

SETUP:
cooper\_pair       : \langleD\_mol; T\_⋈; $R_{\supseteq}$; P\_sym; $F_{\hbar}$; K\_slow; G\_ג; $\Gamma$\_$\land$(SPEC); $\Phi$\_c; $\Omega$\_Z\rangle
conducting\_electron: \langleD\_mol; T\_\textbar{}; $R_{\supseteq}$; P\_donor; $F_{\hbar}$; K\_fast; G\_ב; $\Gamma$\_$\land$(SEL)\rangle

QUESTION: tensor(electron, ?) $\approx$ cooper\_pair
Find the ''pairing partner'' --- the quasiparticle that, when tensored with
the conducting electron, produces the Cooper pair primitive tuple.

COFACTOR RULES (inverting tensor per axis):
F (min-dominant): tensor{[}F{]} = min(A\_F, B\_F)
A\_F = $F_{\hbar}$, C\_F = $F_{\hbar}$ $\rightarrow$ A explains F; B\_F = $F_{\hbar}$ (must be equally high)
K (min-dominant): tensor{[}K{]} = min(A\_K, B\_K)
A\_K = K\_fast, C\_K = K\_slow $\rightarrow$ A is NOT the bottleneck; B\_K = K\_slow (B is the slow partner)
G (max-dominant): tensor{[}G{]} = max(A\_G, B\_G)
A\_G = G\_local, C\_G = G\_meso $\rightarrow$ B\_G = G\_meso (B contributes the coherence length)
T (promotion): C\_T = T\_⋈, A\_T = T\_\textbar{} $\rightarrow$ B must contribute T\_⋈ (B has the pairing loop topology)
$\Phi$ (join-dominant): C\_$\Phi$ = $\Phi$\_c, A\_$\Phi$ = $\Phi$\_sub $\rightarrow$ B\_$\Phi$ = $\Phi$\_c (B carries the criticality)
$\Omega$ (join-dominant): C\_$\Omega$ = $\Omega$\_Z, A\_$\Omega$ = $\Omega$\_trivial $\rightarrow$ B\_$\Omega$ = $\Omega$\_Z (B carries the topological invariant)

PHYSICAL MEANING:
The inferred partner B has: K\_slow (the condensate timescale, not the Fermi velocity),
G\_mesoscale (the coherence length), T\_bowtie (the pairing loop),
$\Phi$\_c (the superfluid phase transition), $\Omega$\_Z (the winding number).
This is the PHONON-DRESSED ELECTRON --- the retarded interaction that
makes the other electron ''look slow'' through phonon exchange.
The cofactor correctly reconstructs the effective retarded partner.

STATUS   : PASS
RESULT   : \langleD\_wedge; T\_bowtie; R\_superset; P\_pm\_sym; F\_hbar; K\_slow; G\_gimel; Gamma\_and(SPECIFIC); Phi\_c; Omega\_Z\rangle

PER-AXIS ROLES:
F    : BOTTLENECK       A is fidelity bottleneck; B $\geq$ Fidelity.HIGH
K    : BOTTLENECK       B sets K floor = KineticCharacter.SLOW
G    : CONTRIBUTOR      B is the G=Granularity.MESOSCALE contributor
D    : EXPLAINED        A explains D; B needs only D\_MOLECULAR
Phi  : CONTRIBUTOR      B must carry $\Phi$\_c (A doesn't have it)
Omega: CONTRIBUTOR      B carries $\Omega$=TopoIndex.Z\_CLASS
T    : PASSTHROUGH      B contributes T=Topology.CYCLIC\_BOWTIE (A has Topology.LINEAR)
R    : EXPLAINED        A explains R=RecognitionMode.NON\_COVALENT
P    : PASSTHROUGH      B contributes P=Polarity.SELF\_COMPLEMENTARY\_SYM (A has Polarity.DONOR)
Gamma: PASSTHROUGH      B contributes Gamma=InteractionGrammar.SPECIFIC\_AND (A has InteractionGrammar.SELECTIVE\_AND)

$\Phi$\_c SOURCE: cofactor
⟹  Cofactor(cooper\_pair \textbar{} conducting\_electron) computed successfully
⟹    F: BOTTLENECK --- A is fidelity bottleneck; B $\geq$ Fidelity.HIGH
⟹    K: BOTTLENECK --- B sets K floor = KineticCharacter.SLOW
⟹    G: CONTRIBUTOR --- B is the G=Granularity.MESOSCALE contributor
⟹    D: EXPLAINED --- A explains D; B needs only D\_MOLECULAR
⟹    Phi: CONTRIBUTOR --- B must carry $\Phi$\_c (A doesn't have it)
⟹    Omega: CONTRIBUTOR --- B carries $\Omega$=TopoIndex.Z\_CLASS
⟹    T: PASSTHROUGH --- B contributes T=Topology.CYCLIC\_BOWTIE (A has Topology.LINEAR)
⟹    R: EXPLAINED --- A explains R=RecognitionMode.NON\_COVALENT
⟹    P: PASSTHROUGH --- B contributes P=Polarity.SELF\_COMPLEMENTARY\_SYM (A has Polarity.DONOR)
⟹    Gamma: PASSTHROUGH --- B contributes Gamma=InteractionGrammar.SPECIFIC\_AND (A has InteractionGrammar.SELECTIVE\_AND)

────────────────────────────────────────────────────────────────────────
7.2  principal\_decomp(GNF2) --- join-irreducible basis decomposition (SVD analog)
────────────────────────────────────────────────────────────────────────

SETUP:
GNF-2 : \langleD\_$\land$; T\_branched; $R_{\supseteq}$; P\_+-; F\_ℇ; K\_mod; G\_ג; $\Gamma$\_$\rightarrow$(SEL); $\Phi$\_c; $\Omega$\_0\rangle

The decomposition produces an ordered list of join-irreducible factors.
Each factor = single primitive contribution above the constraint-bottom.
Reading order = most constraining $\rightarrow$ least constraining.

TENSOR-MATH ANALOGY:
In a product lattice L = L\_F $\times$ L\_K $\times$ L\_G $\times$ L\_categorical,
every element s has a unique Birkhoff representation as a join
of join-irreducible elements. This is the lattice-theoretic analog
of expressing a vector in terms of basis vectors.
The ''principal'' decomposition orders these by their \xi\_CP contribution ---
the analog of ordering singular values by magnitude.

EXPECTED FACTORS (rough order):
1. $\Phi$\_c component   --- allosteric criticality (hardest to satisfy)
2. G\_meso component --- mesoscale propagation
3. K\_mod component  --- kinetic character
4. F\_med component  --- fidelity floor
5. Categorical skeleton (D, T, R, P, $\Gamma$ unchanged)

DESIGN USE: tells you which primitive of GNF-2 is the HARDEST to engineer.
If you're trying to improve GNF-2, start with factor 1 (highest \xi contribution).

FACTORS (4 total, each is a join-irreducible atom):
{[}1{]} atom{[}F=F\_eth{]}                   non-bottom: F=F\_eth
{[}2{]} atom{[}K=K\_mod{]}                   non-bottom: K=K\_mod
{[}3{]} atom{[}G=G\_gimel{]}                 non-bottom: G=G\_gimel
{[}4{]} skeleton(GNF2\_allosteric)       non-bottom: $\Phi$\_c

\xi BALANCE: 0.000 nats
⟹    Factor 1: F contribution = Fidelity.MEDIUM
⟹    Factor 2: K contribution = KineticCharacter.MODERATE
⟹    Factor 3: G contribution = Granularity.MESOSCALE
⟹    Skeleton: categorical primitives of GNF2\_allosteric

────────────────────────────────────────────────────────────────────────
7.3  project + complement\_rel --- Heyting pseudocomplement (complementary design)
────────────────────────────────────────────────────────────────────────

STEP 1: project(cooper\_pair, {[}''criticality\_phase'', ''topo\_index''{]})
Retain only $\Phi$ and $\Omega$; zero out all other primitives.
Analogous to: $\pi$ᵢ(v) --- project vector v onto the \{i\}-th coordinate subspace.

PROJECTED: \langleD\_wedge; T\_linear; R\_superset; P\_pm\_sym; F\_ell; K\_fast; G\_beth; Gamma\_or(BROAD); Phi\_c; Omega\_Z\rangle
RETAINED : $\Phi$ = CriticalityPhase.CRITICAL, $\Omega$ = TopoIndex.Z\_CLASS
ZEROED   : D, T, R, P, F, K, G, Gamma

STEP 2: complement\_rel(gnf2, context=projected, target=cooper\_pair)
Find the maximal x $\leq$ GNF-2 such that:
(1) x ⊓ projected = \perp   (x has no overlap with the $\Phi$/$\Omega$ projection)
(2) x ⊔ projected $\geq$ cooper\_pair  (together they cover the target)

TENSOR-MATH ANALOGY:
In a Heyting algebra (the algebraic model of intuitionistic logic):
a $\Rightarrow$ b  =  ⋁\{x : x $\land$ a $\leq$ b\}   (the IMPLICATION / pseudocomplement)
Here: complement\_rel(GNF-2, context, target) = GNF-2's ''contribution''
to reaching the cooper\_pair target AFTER accounting for what context already covers.

\begin{lstlisting}
This is also the constructive analog of the QUOTIENT in module theory:
GNF-2 / context = the part of GNF-2 that context does NOT already explain.
\end{lstlisting}

DESIGN MEANING:
The result tells you: ''if the $\Phi$/$\Omega$ structure is already given by the
topological material, what does the molecular GNF-2 uniquely contribute
toward the cooper\_pair target?''
Answer: K\_slow + G\_meso + T\_branched (the slow, mesoscale, branched kinetics
that the topological material alone cannot provide).

SATISFIED : False
⟹    K: complement value = KineticCharacter.MBL
⟹    G: complement value = Granularity.MESOSCALE
⟹  F: cannot satisfy target F=Fidelity.HIGH
⟹    F: complement value = Fidelity.MEDIUM
⟹    D: context covers target; x{[}D{]} = Dimensionality.MOLECULAR
⟹    T: x{[}T{]} = Topology.CYCLIC\_BOWTIE (complement of context Topology.LINEAR)
⟹    R: context covers target; x{[}R{]} = RecognitionMode.NON\_COVALENT
⟹    P: context covers target; x{[}P{]} = Polarity.DONOR\_ACCEPTOR
⟹    Gamma: x{[}Gamma{]} = InteractionGrammar.SPECIFIC\_AND (complement of context InteractionGrammar.BROAD\_OR)
⟹    Phi: context covers target; x{[}Phi{]} = None
⟹    Omega: context covers target; x{[}Omega{]} = TopoIndex.TRIVIAL

━━━━━━━━━━━━━━━━━━━━━━━━━━━━━━━━━━━━━━━━━━━━━━━━━━━━━━━━━━━━━━━━━━━━━━━━
Run individual sections:  python TENSOR\_OPS\_DEMO.py --section tensor
━━━━━━━━━━━━━━━━━━━━━━━━━━━━━━━━━━━━━━━━━━━━━━━━━━━━━━━━━━━━━━━━━━━━━━━━


\end{document}
